% LessAI 综合测试（TeX）
% 生成时间：2026-04-09
% 目标：覆盖标题、列表、链接、强调、代码、数学、表格、参考文献和写回敏感区。

\documentclass[a4paper,11pt]{article}
\usepackage{ctex}
\usepackage{amsmath,amssymb}
\usepackage{geometry}
\usepackage[colorlinks,linkcolor=blue]{hyperref}
\usepackage{listings}
\usepackage{longtable}

\geometry{margin=2.4cm}
\linespread{1.2}

\lstset{
  basicstyle=\ttfamily\small,
  breaklines=true,
  frame=single
}

\title{LessAI 综合测试（TeX）}
\author{示例作者}
\date{2026-04-09}

\begin{document}
\maketitle
\tableofcontents

\section{文档说明与测试目标}

这份文档用于验证 TeX 源码在导入后是否还能保持可读结构、稳定分块和安全写回。
建议分别切换段落、整句、小句三种粒度，并观察自然语言区域和命令保护区的边界是否清晰。

TeX 的关键点在于：命令、环境、数学、代码和注释不能被误改；正文中的自然语言又不能被这些结构切得支离破碎。

\section{标题层级与普通段落}

\subsection{小标题示例}

这一段是普通正文。它包含中文句号、问号和感叹号，也包含小数 3.14、日期 2026-04-09、版本号 v2.3.7。整句模式下，句末应成为主要边界；段落模式下，这一整段应保持连续。

下一段继续模拟报告正文：先描述背景，再给出判断，最后补一句限制条件。这里的目标不是考验排版，而是观察长句在不同粒度下是否仍然可读。

\section{标点边界与数字边界}

这一段包含中英混排：Dr.\ Wang 在 U.S.A.\ 完成实验后，于 2025-12-31 提交了 draft v1.0.2。若整句模式对缩写处理不够理想，可以记录问题，但段落级阅读不应被破坏。

还要观察括号、引号和书名号的连续性：例如“引用内容”、（局部说明）、《示例报告》。这些成对符号不应被错误吞掉。

\section{列表与段内换行}

下面先给一个显式段落，再用 \texttt{\textbackslash par} 制造真正的段落边界。
\par
这一段用于验证：段落级分块应更接近渲染后的段落，而不是逐行切碎。

\begin{enumerate}
  \item 列表项 1：包含两个整句。第一句结束。第二句继续，用于观察整句边界。
  \item 列表项 2：包含逗号、分号；用于观察小句边界：例如，我们希望在“，”“；”“：”附近切开，但不要改坏骨架。
  \item 列表项 3：包含一个换行命令：\\[0.5em]
  它不应被误识别为数学块的开始。
\end{enumerate}

\begin{itemize}
  \item 项目 A：包含路径 \texttt{E:\textbackslash Code\textbackslash LessAI\textbackslash testdoc}。
  \item 项目 B：包含命令 token \texttt{cargo fmt --check}。
  \item 项目 C：包含一个链接 \href{https://example.com/list}{https://example.com/list}。
\end{itemize}

\section{链接与可读正文}

这一段集中放常见链接命令：\href{https://example.com/docs}{官方说明页}、\url{https://example.com/url-only}、\path|E:/Code/LessAI/src-tauri|。这些结构应该被锁定，但不应打断上下文阅读。

如果你只改动这段里的自然语言，例如把“集中放”改成“专门放”，理想效果是链接命令仍保持原样，局部 diff 也只发生在正文区域。

\section{行内强调与常见样式}

这里放常见文本命令：\textbf{加粗内容}、\emph{强调内容}、\underline{下划线内容}、\textit{斜体内容}、\texttt{monospace}，以及脚注\footnote{脚注正文通常也是文章的一部分，但命令骨架不应被改坏。}。

这些命令的语法骨架需要保留，命令包裹的自然语言应尽量维持正常阅读和稳定分块。

\section{代码、命令与原样内容}

本节先测试行内原样命令：\verb|pnpm run tauri:dev|、\lstinline|cargo test docx -- --nocapture|。

\begin{lstlisting}
pnpm install
pnpm run tauri:dev
cargo test
echo "a%b"
\end{lstlisting}

\begin{verbatim}
if score >= 90:
    level = "A"
else:
    level = "B"
\end{verbatim}

代码环境应整体保护，环境前后的正文则应恢复正常分块。

\section{数学与公式相关内容}

这一段包含行内公式 $E=mc^2$、$\alpha+\beta\to\gamma$，以及另一种写法 \(\bar{x}=\frac{1}{n}\sum_{i=1}^{n}x_i\)。这些公式不应被正文改写破坏。

\[
  \hat{\mu} = \frac{1}{n}\sum_{i=1}^{n} x_i,\qquad
  \hat{\Sigma} = \frac{1}{n-1}\sum_{i=1}^{n}(x_i-\hat{\mu})(x_i-\hat{\mu})^\top
\]

\begin{equation}
  p = 2\left[1-\Phi(2.6980)\right]
\end{equation}

\begin{align}
  \bar{x} &= \frac{1}{n}\sum_{i=1}^{n} x_i, \\
  s^2 &= \frac{1}{n-1}\sum_{i=1}^{n}(x_i-\bar{x})^2
\end{align}

公式前后的自然语言仍应能连续阅读，而不是被切成毫无关联的小碎块。

\section{表格与非正文结构}

下面是一个 \texttt{tabular}，它通常应作为结构性强的保护区保留原样：
\begin{center}
\begin{tabular}{l|c|c}
项目 & 期望行为 & 备注 \\
\hline
标题 & 保留命令骨架 & 不要误删反斜杠 \\
链接 & 保留目标与显示文本 & 不要改坏 URL \\
数学 & 保留公式边界 & 前后正文可读 \\
\end{tabular}
\end{center}

\begin{thebibliography}{9}
\bibitem{ref1} Example Author. \emph{An Example Book}. 2026.
\bibitem{ref2} Another Author. ``A short paper title.'' 2025.
\end{thebibliography}

\section{混合段落与锁定区穿插}

这一段故意把多种结构放在一起：先写普通正文，再插入 \href{https://example.com/mixed}{链接文本}，再插入 \texttt{git status}、\verb|pnpm run build|、$a_i=\frac{1}{n}$，最后补一句“这些保护区不应打断整体可读性”。

% 这是一行注释：它不应把同一段误切成新块
这一句与上一句之间只有注释，没有真正的段落边界，因此段落模式下仍应视为同一段。

\smallskip
这一句前面只有视觉空白命令，不应自动变成新的段落块。

\section{写回敏感区}

请只做下面一种最小修改：把某个逗号改成分号，或者在两句之间插入一个空行，或者把问号改成感叹号。

短句 A，用来观察局部 diff。
短句 B，用来观察空行边界。
短句 C？

\par

如果写回后重新导入，短句顺序、空行位置和相邻段落边界都应保持稳定。

\section{长段落与压力观察区}

这一整段故意写得比较长，用来观察 TeX 正文在复杂混排下的稳定性：当一段文字同时包含中文叙述、English phrases、缩写 e.g.\ 和 i.e.、数字 0.618、日期 2024-2026、链接 \url{https://example.com/report/final}、命令 \texttt{cargo test}、路径 \path|/mnt/e/code/lessai/testdoc|、公式 $f(x)=x^2$、以及括号（例如这里的补充说明）时，段落模式应优先保证整体可读，整句模式应主要在真正句末切开，小句模式则可以在逗号、分号、冒号附近进一步细分，但任何一种模式都不应遗漏命令、误识别命令，或把一个自然段打成一串无意义的小块。

\end{document}
